\small
\begin{enumerate}[label=\textbf{\arabic*.}, leftmargin=0pt, itemsep=1.2em]
\item Как обычно представляют решение задачи коммивояжера в генетическом алгоритме?
\begin{itemize}[label=$\square$, leftmargin=1.5cm, itemsep=0.2em, topsep=0.3em]
\item Как перестановку индексов городов
\item Как бинарную строку, где каждый бит означает наличие города
\item Как матрицу ошибок прогноза
\item Как одно вещественное число
\end{itemize}
{\small\textit{Правильный ответ: Как перестановку индексов городов.}}

\item Что означает допустимый тур в задаче коммивояжера?
\begin{itemize}[label=$\square$, leftmargin=1.5cm, itemsep=0.2em, topsep=0.3em]
\item Каждый город посещается ровно один раз, после чего происходит возврат в исходный город
\item Маршрут обязательно проходит через один и тот же город несколько раз
\item Маршрут заканчивается в произвольной точке без возврата
\item Разрешены пропуски части городов
\end{itemize}
{\small\textit{Правильный ответ: Каждый город посещается ровно один раз, после чего происходит возврат в исходный город.}}

\item Как обычно вычисляется длина тура?
\begin{itemize}[label=$\square$, leftmargin=1.5cm, itemsep=0.2em, topsep=0.3em]
\item Как сумма расстояний между последовательными городами с учетом возврата в начало
\item Как произведение всех расстояний между городами
\item Как число городов в маршруте
\item Как среднее значение координат городов
\end{itemize}
{\small\textit{Правильный ответ: Как сумма расстояний между последовательными городами с учетом возврата в начало.}}

\item Почему для задачи коммивояжера нельзя напрямую использовать обычный одноточечный кроссинговер бинарного ГА?
\begin{itemize}[label=$\square$, leftmargin=1.5cm, itemsep=0.2em, topsep=0.3em]
\item Потому что он легко порождает недопустимые перестановки с повторами и пропусками
\item Потому что кроссинговер в TSP вообще не нужен
\item Потому что у TSP нет родителей и потомков
\item Потому что длина тура не зависит от порядка городов
\end{itemize}
{\small\textit{Правильный ответ: Потому что он легко порождает недопустимые перестановки с повторами и пропусками.}}

\item В чем идея циклического кроссинговера (CX)?
\begin{itemize}[label=$\square$, leftmargin=1.5cm, itemsep=0.2em, topsep=0.3em]
\item Он сохраняет циклы соответствия позиций элементов между двумя родителями
\item Он случайно разворачивает сегмент маршрута
\item Он всегда усредняет индексы городов
\item Он выбирает только первый город и удаляет остальные
\end{itemize}
{\small\textit{Правильный ответ: Он сохраняет циклы соответствия позиций элементов между двумя родителями.}}

\item Какое важное свойство потомков обеспечивает CX-кроссинговер?
\begin{itemize}[label=$\square$, leftmargin=1.5cm, itemsep=0.2em, topsep=0.3em]
\item Потомки остаются допустимыми перестановками без повторов и пропусков
\item Потомки всегда короче обоих родителей
\item Потомки содержат только города первого родителя
\item Потомки всегда совпадают с оптимальным маршрутом
\end{itemize}
{\small\textit{Правильный ответ: Потомки остаются допустимыми перестановками без повторов и пропусков.}}

\item Что делает swap mutation в TSP?
\begin{itemize}[label=$\square$, leftmargin=1.5cm, itemsep=0.2em, topsep=0.3em]
\item Меняет местами два города в перестановке
\item Удаляет из маршрута два города
\item Разворачивает весь маршрут полностью
\item Заменяет координаты городов случайным образом
\end{itemize}
{\small\textit{Правильный ответ: Меняет местами два города в перестановке.}}

\item Почему swap mutation удобна для TSP?
\begin{itemize}[label=$\square$, leftmargin=1.5cm, itemsep=0.2em, topsep=0.3em]
\item Она сохраняет корректность перестановки
\item Она всегда улучшает маршрут
\item Она устраняет необходимость в селекции
\item Она делает длину тура постоянной
\end{itemize}
{\small\textit{Правильный ответ: Она сохраняет корректность перестановки.}}

\item Что означает элитарная стратегия в генетическом алгоритме для TSP?
\begin{itemize}[label=$\square$, leftmargin=1.5cm, itemsep=0.2em, topsep=0.3em]
\item Лучшие туры переходят в следующее поколение без ухудшения
\item Сохраняются только случайные туры
\item Каждый потомок обязательно заменяет лучшего родителя
\item Селекция не используется
\end{itemize}
{\small\textit{Правильный ответ: Лучшие туры переходят в следующее поколение без ухудшения.}}

\item Как работает турнирная селекция применительно к TSP?
\begin{itemize}[label=$\square$, leftmargin=1.5cm, itemsep=0.2em, topsep=0.3em]
\item Из нескольких случайных маршрутов выбирается маршрут с наименьшей длиной
\item Из нескольких случайных маршрутов выбирается самый длинный
\item Родитель определяется только по номеру поколения
\item Маршруты выбираются без учета длины
\end{itemize}
{\small\textit{Правильный ответ: Из нескольких случайных маршрутов выбирается маршрут с наименьшей длиной.}}

\item Почему задача коммивояжера относится к задачам комбинаторной оптимизации?
\begin{itemize}[label=$\square$, leftmargin=1.5cm, itemsep=0.2em, topsep=0.3em]
\item Потому что требуется выбрать наилучшую перестановку из огромного конечного множества маршрутов
\item Потому что ее решение всегда выражается одной формулой
\item Потому что координаты городов непрерывно изменяются во времени
\item Потому что задача сводится к линейной регрессии
\end{itemize}
{\small\textit{Правильный ответ: Потому что требуется выбрать наилучшую перестановку из огромного конечного множества маршрутов.}}

\item Как растет число возможных маршрутов при увеличении числа городов?
\begin{itemize}[label=$\square$, leftmargin=1.5cm, itemsep=0.2em, topsep=0.3em]
\item Очень быстро, факториально
\item Линейно
\item Квадратично
\item Оно почти не зависит от числа городов
\end{itemize}
{\small\textit{Правильный ответ: Очень быстро, факториально.}}

\item Почему TSP считается сложной задачей для полного перебора?
\begin{itemize}[label=$\square$, leftmargin=1.5cm, itemsep=0.2em, topsep=0.3em]
\item Число допустимых туров слишком велико даже для умеренного числа городов
\item Потому что невозможно вычислить расстояния между городами
\item Потому что маршрут всегда единственный
\item Потому что координаты городов неизвестны
\end{itemize}
{\small\textit{Правильный ответ: Число допустимых туров слишком велико даже для умеренного числа городов.}}

\item Для чего используют тестовые наборы TSPLIB, такие как eil76?
\begin{itemize}[label=$\square$, leftmargin=1.5cm, itemsep=0.2em, topsep=0.3em]
\item Для объективного сравнения алгоритмов на стандартных экземплярах задач
\item Только для построения нейронных сетей Кохонена
\item Только для задач классификации
\item Только для визуализации графиков потерь
\end{itemize}
{\small\textit{Правильный ответ: Для объективного сравнения алгоритмов на стандартных экземплярах задач.}}

\item Что означает известный оптимум 538 для eil76?
\begin{itemize}[label=$\square$, leftmargin=1.5cm, itemsep=0.2em, topsep=0.3em]
\item Это длина лучшего известного тура для данного набора
\item Это число городов в задаче
\item Это размер популяции по умолчанию
\item Это номер поколения, на котором алгоритм останавливается
\end{itemize}
{\small\textit{Правильный ответ: Это длина лучшего известного тура для данного набора.}}

\item Какую роль играет вероятность мутации в TSP-алгоритме?
\begin{itemize}[label=$\square$, leftmargin=1.5cm, itemsep=0.2em, topsep=0.3em]
\item Помогает выходить из локальных оптимумов за счет перестройки маршрута
\item Полностью заменяет кроссинговер
\item Определяет число городов в маршруте
\item Не влияет на поиск
\end{itemize}
{\small\textit{Правильный ответ: Помогает выходить из локальных оптимумов за счет перестройки маршрута.}}

\item Что может произойти при слишком малой вероятности мутации в TSP?
\begin{itemize}[label=$\square$, leftmargin=1.5cm, itemsep=0.2em, topsep=0.3em]
\item Популяция может быстро стать однообразной и застрять в локальном оптимуме
\item Маршрут станет недопустимым даже без операторов
\item Алгоритм перестанет считать длину тура
\item Все туры станут оптимальными
\end{itemize}
{\small\textit{Правильный ответ: Популяция может быстро стать однообразной и застрять в локальном оптимуме.}}

\item Что может произойти при слишком высокой вероятности мутации в TSP?
\begin{itemize}[label=$\square$, leftmargin=1.5cm, itemsep=0.2em, topsep=0.3em]
\item Хорошие маршруты будут разрушаться слишком часто
\item Поиск станет детерминированным и стабильным
\item Кроссинговер станет не нужен математически
\item Оптимальный тур будет найден гарантированно
\end{itemize}
{\small\textit{Правильный ответ: Хорошие маршруты будут разрушаться слишком часто.}}

\item Зачем визуализируют лучший найденный тур на карте городов?
\begin{itemize}[label=$\square$, leftmargin=1.5cm, itemsep=0.2em, topsep=0.3em]
\item Чтобы наглядно оценить структуру маршрута и качество решения
\item Чтобы заменить вычисление длины тура
\item Чтобы вычислить коэффициент силуэта
\item Чтобы выбрать размер скрытого слоя
\end{itemize}
{\small\textit{Правильный ответ: Чтобы наглядно оценить структуру маршрута и качество решения.}}

\item Что обычно показывают графики сходимости в TSP-экспериментах?
\begin{itemize}[label=$\square$, leftmargin=1.5cm, itemsep=0.2em, topsep=0.3em]
\item Как меняется длина лучшего и среднего маршрута по поколениям
\item Как изменяется число нейронов в сети
\item Как изменяется число признаков в данных
\item Как изменяется значение learning rate в градиентном спуске
\end{itemize}
{\small\textit{Правильный ответ: Как меняется длина лучшего и среднего маршрута по поколениям.}}

\item Почему для TSP важны специальные операторы для перестановок?
\begin{itemize}[label=$\square$, leftmargin=1.5cm, itemsep=0.2em, topsep=0.3em]
\item Потому что решение должно оставаться корректной последовательностью городов
\item Потому что обычные вещественные операторы всегда быстрее
\item Потому что длина тура не зависит от порядка
\item Потому что TSP не использует популяции
\end{itemize}
{\small\textit{Правильный ответ: Потому что решение должно оставаться корректной последовательностью городов.}}

\item Что означает локальный оптимум в задаче коммивояжера?
\begin{itemize}[label=$\square$, leftmargin=1.5cm, itemsep=0.2em, topsep=0.3em]
\item Маршрут, который нельзя заметно улучшить небольшими изменениями, но который не обязательно глобально лучший
\item Маршрут с нулевой длиной
\item Маршрут без возврата в исходный город
\item Маршрут, содержащий только один город
\end{itemize}
{\small\textit{Правильный ответ: Маршрут, который нельзя заметно улучшить небольшими изменениями, но который не обязательно глобально лучший.}}

\item Почему эволюционный подход полезен для TSP?
\begin{itemize}[label=$\square$, leftmargin=1.5cm, itemsep=0.2em, topsep=0.3em]
\item Он позволяет исследовать большое пространство маршрутов без полного перебора
\item Он делает задачу выпуклой
\item Он не требует вычисления расстояний
\item Он гарантирует нахождение точного оптимума за фиксированное число поколений
\end{itemize}
{\small\textit{Правильный ответ: Он позволяет исследовать большое пространство маршрутов без полного перебора.}}

\item Что можно сказать о качестве решения ГА для TSP?
\begin{itemize}[label=$\square$, leftmargin=1.5cm, itemsep=0.2em, topsep=0.3em]
\item Обычно алгоритм находит хорошие субоптимальные туры, но не гарантирует точный оптимум
\item Алгоритм всегда дает точный оптимум
\item Алгоритм не способен улучшить случайный маршрут
\item Качество решения не зависит от операторов
\end{itemize}
{\small\textit{Правильный ответ: Обычно алгоритм находит хорошие субоптимальные туры, но не гарантирует точный оптимум.}}

\item Какой эффект часто дает увеличение размера популяции в TSP?
\begin{itemize}[label=$\square$, leftmargin=1.5cm, itemsep=0.2em, topsep=0.3em]
\item Повышает разнообразие туров и улучшает шансы найти качественные решения
\item Делает туры заведомо короче оптимума
\item Отменяет необходимость мутации
\item Превращает перестановку в бинарную строку
\end{itemize}
{\small\textit{Правильный ответ: Повышает разнообразие туров и улучшает шансы найти качественные решения.}}

\item Зачем в TSP иногда комбинируют ГА с локальным поиском вроде 2-opt?
\begin{itemize}[label=$\square$, leftmargin=1.5cm, itemsep=0.2em, topsep=0.3em]
\item Чтобы дополнительно улучшать уже найденные хорошие маршруты
\item Чтобы заменить функцию длины тура
\item Чтобы исключить селекцию
\item Чтобы перевести задачу в регрессию
\end{itemize}
{\small\textit{Правильный ответ: Чтобы дополнительно улучшать уже найденные хорошие маршруты.}}

\item Какое утверждение о CX-кроссинговере наиболее корректно?
\begin{itemize}[label=$\square$, leftmargin=1.5cm, itemsep=0.2em, topsep=0.3em]
\item Он ориентирован на позиционную информацию в перестановках
\item Он работает только для бинарных строк
\item Он всегда возвращает одного потомка
\item Он не зависит от порядка элементов в родителях
\end{itemize}
{\small\textit{Правильный ответ: Он ориентирован на позиционную информацию в перестановках.}}

\item Что является главной целью фитнес-функции в TSP?
\begin{itemize}[label=$\square$, leftmargin=1.5cm, itemsep=0.2em, topsep=0.3em]
\item Оценить качество маршрута, обычно через его длину
\item Определить число поколений до остановки
\item Сгенерировать новых родителей
\item Нормализовать координаты городов
\end{itemize}
{\small\textit{Правильный ответ: Оценить качество маршрута, обычно через его длину.}}

\item Какое условие обязательно для корректного потомка в TSP?
\begin{itemize}[label=$\square$, leftmargin=1.5cm, itemsep=0.2em, topsep=0.3em]
\item Каждый город должен встречаться ровно один раз
\item Один и тот же город должен появляться минимум дважды
\item Порядок городов не играет роли
\item Маршрут может быть незамкнутым
\end{itemize}
{\small\textit{Правильный ответ: Каждый город должен встречаться ровно один раз.}}

\item Какой из факторов особенно важен для успешного ГА на TSP?
\begin{itemize}[label=$\square$, leftmargin=1.5cm, itemsep=0.2em, topsep=0.3em]
\item Подходящее сочетание представления, операторов перестановок и параметров поиска
\item Использование только линейной функции активации
\item Наличие обучающей выборки с метками классов
\item Отказ от оценки длины маршрута
\end{itemize}
{\small\textit{Правильный ответ: Подходящее сочетание представления, операторов перестановок и параметров поиска.}}

\end{enumerate}
